\section{Discussion}


This section organizes the positioning of this paper, differences from existing theories, theoretical implications, and future tasks.

The central claim of this paper is that consciousness should not be treated as a product.

Consciousness is not modeled here as an object produced somewhere.

Consciousness is a log-referencing state that opens when a certain speed difference, coupling viscosity, internal delay, and meaning gradient obtain between fast predictive processes and slow log-sedimentation processes.

This view does not erase consciousness.

Rather, it removes consciousness from its conventional positions — as source of causation, object within the brain, content of representation, container of qualia — and repositions it as a temporal non-closure regime.

\subsection{Relation to Global Workspace-Type Models}


In Global Workspace-type theories, consciousness is often understood as information becoming globally available, or being broadcast to multiple processing systems.

This paper can connect with this orientation in part.

In particular, the point that conscious states obtain not at a single site but as a relation among multiple processing systems has affinity with the idea of global availability.

But this paper does not define consciousness as "broadcasting" or "sharing" itself.

What this paper emphasizes, before information is widely used, is in which time scale it becomes referenceable.

Fast-layer processing sediments into the slow layer and becomes referenceable as a log.

At that point, the conscious state opens.

Spatial breadth is therefore less important in this paper than temporal transformation.

Information becoming broadly available may be part of the result or conditions of conscious state, but it is not consciousness itself.

\subsection{Relation to IIT-Type Models}


In IIT-type theories, consciousness is associated with the degree of information integration or structural indivisibility.

This paper also does not reduce consciousness to simple local processing, and emphasizes that multiple elements form a relational structure.

In that sense, treating consciousness as structural relation rather than isolated parts has some common ground.

But this paper does not define consciousness as an integration quantity.

Consciousness is treated not as the magnitude of integrated information but as the temporal condition under which log referencing opens.

If integration is too strong, structure becomes fixed and may lose plasticity.

If integration is too weak, processing dissipates and log referencing cannot obtain.

In this paper, therefore, the intermediate state that does not reach complete closure is emphasized over integration itself.

Consciousness is not maximum integration but a non-closure stability window.

\subsection{Relation to Higher-Order and Self-Model Theories}


In Higher-Order and self-model theories, consciousness is explained as one representation capturing one's own state, or as a self-model representing a state.

This paper acknowledges the importance of self-reference and re-referencing.

Indeed, for log referencing to obtain, current processing must have a relation to slow-layer logs, and that relation must in turn influence current processing.

In this sense, this paper does not deny self-reference.

But this paper does not hold that self-representation generates consciousness.

The self and ego are stable structures of the slow layer formed through the repetition of log referencing, and are not the cause of consciousness.

The ego is not a subject but the position at which speed difference is transformed.

Self-models are important, but they are terrain formed within the log-referencing loop, not the source of consciousness generation.

\subsection{Relation to Predictive Processing and Free Energy Approaches}


This paper emphasizes prediction, error, and forward models in the fast layer.

In this respect, it can connect with predictive processing and the free energy principle.

The fast layer anticipates the future, generating action candidates and error corrections.

The slow layer sediments part of this and changes the terrain for subsequent prediction.

This loop contains the structure of prediction and updating.

But the focus of this paper is not the accuracy of prediction or minimization of free energy itself.

What matters is how prediction sediments into a log-referenceable structure.

Even with fast prediction, if it does not reach the slow layer, the conscious state does not open.

Conversely, if sedimentation is too strong, prediction becomes fixed to specific meaning wells.

This paper therefore does not view predictive processing as a sufficient condition for consciousness.

Prediction may be a necessary element, but conscious states require speed difference, viscosity, delay, meaning gradient, and non-closure.

\subsection{Why the Model Is Not Reductionist}


This paper does not treat consciousness as a mysterious entity.

But this does not mean reducing consciousness to simple physical quantities or neural activity.

Rather, to avoid simple reduction, this paper treats consciousness as a state domain.

Consciousness is not $I$ itself, not $J$, not $\Phi$, not any one of $\Delta v$, $\mu$, $\tau$, $\mathcal{T}_{\Phi}$.

Consciousness is the regime that opens when these enter a certain relation.

In this sense, the position of this paper is neither eliminativist nor dualist.

It does not add consciousness as a special entity.

But it does not collapse consciousness into a single physical quantity either.

It treats consciousness as the relation among multiple layers.

\subsection{Why the Model Is Not Anti-Neuroscience}


This paper holds that consciousness does not localize to a single site.

But this is not a negation of neuroscience.

Rather, it is a framework for interpreting neuroscientific findings more carefully.

Certain brain regions, circuits, signals, oscillations, and activity patterns can certainly correlate with conscious states.

These are important local substrates.

But those local substrates cannot be identified with consciousness itself.

In the position of this paper, neural activity is interpreted as part of the log-referencing loop.

Which activity corresponds to the fast layer? Which activity corresponds to sedimentation into the slow layer? Which pathway supports reportability? Which activity forms the meaning gradient?

Asking in this way repositions neuroscientific data not as searching for the seat of consciousness but as investigating the conditions under which conscious states open.

\subsection{The Role of IFGT and VED}


In this paper, IFGT and VED have been used not as completed theories that fully explain consciousness but as a coordinate framework for repositioning consciousness.

IFGT provides the structure by which log referencing is directed along the meaning gradient, through projected uses of information density $I$, information flow $J$, and information potential $\Phi$.

VED provides the structure of difference, flow, and persistence in which complete closure is not realized as an ordinary state.

Connecting these two, consciousness is positioned as follows:

Consciousness is modeled as the state in which log referencing directed by the projected information potential is open within a temporal non-closure regime that reaches neither complete synchrony nor complete separation.

This formulation does not reduce consciousness to a physical field.

It is a coordinate system for fixing the variables and relations required when speaking of consciousness.

\subsection{Theoretical Implications}


The theoretical implications of this paper include several points.

First, the question in consciousness research shifts from "where is consciousness generated?" to "under which conditions does log referencing open?"

Second, qualia are repositioned not as the substance of consciousness but as the live foreground at the sedimentation boundary.

Third, measurement is understood not as capturing consciousness as it is but as an operation converting live openings into sedimented residue.

Fourth, neural correlates of consciousness are interpreted not as consciousness itself but as local conditions participating in the log-referencing loop.

Fifth, pathologies and altered states can be structurally organized not as disease classifications but as directions of deviation from the consciousness window.

These do not negate consciousness research.

Rather, they are translations that move existing data and discussions from the generation model to the opening model.

\subsection{Remaining Questions}


Unresolved tasks remain in this paper.

First, the measurable definition of $\mathcal{T}_{\Phi}$ is incomplete. How to map speed difference, coupling viscosity, and meaning gradient onto neural activity, bodily states, subjective reports, and linguistic logs is a future task.

Second, how to map IFGT's $I$, $J$, $\Phi$ onto brain data and behavioral data must be further refined.

Third, how to empirically estimate the consciousness window remains undetermined.

Fourth, methodological ingenuity is needed for how to experimentally handle qualia sedimentation residue.

Fifth, how the self, language, and social interaction extend the log-referencing structure exceeds the scope of this paper.

These tasks do not negate this paper.

Rather, they show that this paper not only remains unfinished as a theory but also corresponds to the non-closure of consciousness itself as an object.

\subsection{Toward Consciousness, Self, and Intelligence}


This paper treated consciousness as the opening of log referencing.

But this framework does not close with consciousness alone.

The self can be understood as the stable terrain formed in the slow layer through the repetition of openings.

Intelligence can be understood as the structure that uses that terrain to extend future prediction, language, social logs, and interaction with the external environment.

In this sense, this paper is not the endpoint of consciousness theory but the bridge to self theory and intelligence theory.

Consciousness is not modeled as a generated object. Consciousness opens.

The self sediments through the repetition of that opening.

Intelligence flows back into the world again using that sedimented terrain.

Following this flow is the next task.

\subsection{Summary}


This section organized the positioning of this paper.

\begin{itemize}
\item This paper treats consciousness not as a product but as the opening of log referencing.
\item It has points of contact with Global Workspace, IIT, Higher-Order, and Predictive Processing discussions, but is not identical to them.
\item This paper does not reduce consciousness to a single site, single quantity, or single representation.
\item It does not deny neuroscience; it reinterprets neural correlates as local conditions of the log-referencing loop.
\item IFGT and VED are used not as theories that fully explain consciousness but as coordinate systems for positioning it.
\item The claim of this paper is shifting the question in consciousness research from "generation" to "opening conditions."
\item Unresolved tasks remain, but this corresponds not only to the incompleteness of the theory but also to the non-closure of consciousness as its object.
\end{itemize}

The significance of this paper therefore does not lie in explaining consciousness as a closed object.

It lies in repositioning consciousness — as a log-referencing state that opens within temporal non-closure — in a form that allows it to be questioned again.
